\documentclass[11pt]{article}
\usepackage[margin=25mm]{geometry}
\usepackage{fontspec}
\setmainfont{Latin Modern Roman}
\usepackage{amsmath,amssymb,mathtools}
\usepackage{unicode-math}
\setmathfont{Latin Modern Math}
\usepackage{xurl,hyperref}
\hypersetup{colorlinks=true,urlcolor=blue}
\setcounter{secnumdepth}{0}
\setlength{\parindent}{0pt}
\setlength{\parskip}{5pt}
\setlength{\emergencystretch}{3em}
\newcommand{\sref}[2]{\hyperlink{source:#1:#2}{[S#2]}}
\newcommand{\cataloguescope}[1]{\paragraph{Recorded target.}#1}
\begin{document}
\section{The Kissing Number Problem}\label{statement}

\subsection{Definitions and mathematical statement}

\paragraph{Shared notation.}$\mathbb N=\{1,2,\ldots\}$, and $\mathbb Z,\mathbb Q,\mathbb R,\mathbb C$ denote the integers, rationals, reals and complex numbers. Any other conventions are specified locally.


Write $S^{n-1}=\{x\in\mathbb R^n:\lVert x\rVert=1\}$, with the Euclidean inner product $\langle\cdot,\cdot\rangle$. The kissing number is
\[
\tau_n=\max\{|C|:C\subset S^{n-1},\ \langle x,y\rangle\leq\tfrac12\text{ whenever }x\ne y\text{ in }C\}.
\]
The maximum is finite. Multiplying points of $C$ by $2$ gives centers of disjoint-interior unit balls tangent to the unit ball centered at the origin.
\paragraph{Statement recorded in the supplied source.}
Determine $\tau_n$ for every integer $n\geq1$, that is, determine the maximum number of pairwise nonoverlapping unit balls in $\mathbb R^n$ that can simultaneously touch a central unit ball. \sref{statement}{2}

\subsection{Short English statement}
What is the largest possible number of equal spheres that can touch one central sphere without overlapping?
\subsection{Sources}
\begin{description}
\item[\hypertarget{source:statement:1}{S1}] ULAM.AI and the UnsolvedMath Contributors, UnsolvedMath: A Curated Collection of Open Mathematics Problems, catalogue entry. CC BY 4.0. \url{https://huggingface.co/datasets/ulamai/UnsolvedMath}.
\item[\hypertarget{source:statement:2}{S2}] Henry Cohn and Anqi Li, \emph{Improved kissing numbers in seventeen through twenty-one dimensions}, arXiv:2411.04916v3 (2026), Section 1. \url{https://arxiv.org/pdf/2411.04916}.
\end{description}
\label{end:statement}
\end{document}
